\documentclass[10pt]{article}

% ---------- Document Identity (update these only) ----------
\newcommand{\DocStage}{}
\newcommand{\DocTitle}{Your Road to Tier One SOF}
\newcommand{\ProgramName}{182-Week Civilian-to-Selection Readiness Pipeline}
\newcommand{\ProgramSubtitle}{}

% ---------- Page ----------
\usepackage[legalpaper,left=0.5in,right=0.5in,top=0.6in,bottom=0.45in]{geometry}

% ---------- Typography ----------
\usepackage[T1]{fontenc}
\usepackage[utf8]{inputenc}
\usepackage{libertinus}
\usepackage{microtype}
\usepackage{parskip}
\usepackage{ragged2e}
\usepackage{textcomp}
\AtBeginDocument{\color{black}}
\AtBeginDocument{\pagecolor{white}}
\AtBeginDocument{\justifying}

% ---------- Landscape pages ----------
\usepackage{pdflscape}

% ---------- Color ----------
\usepackage[table]{xcolor}
\definecolor{StageBlue}{HTML}{1F4E79}
\definecolor{StageBlueDark}{HTML}{163A5A}
\definecolor{LightGray}{HTML}{F2F4F7}
\definecolor{MidGray}{HTML}{D9DEE5}
\definecolor{RuleGray}{HTML}{C7CED8}
\definecolor{HeaderInk}{HTML}{000000}
\definecolor{Ink}{HTML}{000000}

% ---------- Utility ----------
\usepackage{enumitem}
\sloppy
\setlength{\emergencystretch}{2em}

% ---------- Boxes / UI ----------
\usepackage[most]{tcolorbox}

\tcbset{
  charterTitle/.style={
    enhanced,
    colback=StageBlue!4,
    colframe=StageBlue,
    boxrule=1.25pt,
    arc=3pt,
    left=16pt,right=16pt,top=14pt,bottom=14pt,
    borderline north={3.2pt}{0pt}{StageBlue},
    borderline west={4pt}{0pt}{StageBlue},
    borderline south={1.25pt}{0pt}{StageBlue},
    drop shadow={black!25!white},
  },
  charterRule/.style={
    enhanced,
    colback=LightGray,
    colframe=RuleGray,
    boxrule=0.85pt,
    arc=3pt,
    left=12pt,right=12pt,top=10pt,bottom=10pt,
    borderline west={3pt}{0pt}{StageBlue},
  },
  charterBadge/.style={
    enhanced,
    colback=StageBlue,
    coltext=white,
    colframe=StageBlueDark,
    boxrule=0.6pt,
    arc=2pt,
    left=6pt,right=6pt,top=3pt,bottom=3pt,
  }
}
\newtcbox{\Badge}{nobeforeafter,charterBadge}

% ---------- Lists ----------
\setlist[itemize]{leftmargin=1.5em, itemsep=0.25em, topsep=0.25em}
\setlist[enumerate]{leftmargin=1.8em, itemsep=0.25em, topsep=0.25em}

% ---------- TOC ----------
\usepackage{tocloft}
\setcounter{tocdepth}{3}
\setlength{\cftbeforesecskip}{3pt}
\setlength{\cftbeforesubsecskip}{2pt}
\setlength{\cftbeforesubsubsecskip}{0pt}
\renewcommand{\cftsecfont}{\bfseries}
\renewcommand{\cftsecpagefont}{\bfseries}
\renewcommand{\cftsubsecfont}{\normalfont}
\renewcommand{\cftsubsecpagefont}{\normalfont}
\renewcommand{\cftsubsubsecfont}{\normalfont}
\renewcommand{\cftsubsubsecpagefont}{\normalfont}
\setlength{\cftsecindent}{0pt}
\setlength{\cftsubsecindent}{1.25em}
\setlength{\cftsubsubsecindent}{2.8em}
\setlength{\cftsecnumwidth}{2.1em}
\setlength{\cftsubsecnumwidth}{2.8em}
\setlength{\cftsubsubsecnumwidth}{3.6em}

% Remove the built-in TOC heading so only the custom heading is shown
\makeatletter
\renewcommand{\tableofcontents}{\@starttoc{toc}}
\makeatother

% ---------- Header/Footer: page number only ----------
\usepackage{fancyhdr}
\pagestyle{fancy}
\fancyhf{}
\renewcommand{\headrulewidth}{0pt}
\renewcommand{\footrulewidth}{0pt}
\fancyfoot[C]{\thepage}

% ---------- Section formatting ----------
\usepackage{titlesec}

\titleformat{\section}
  {\Large\bfseries\color{Ink}}
  {\thesection}{0.6em}{}
  [\vspace{0.25em}\textcolor{StageBlue}{\rule{\linewidth}{0.8pt}}]

\titleformat{\subsection}
  {\large\bfseries\color{Ink}}
  {\thesubsection}{0.6em}{}

\titleformat{\subsubsection}
  {\normalsize\bfseries\color{Ink}}
  {\thesubsubsection}{0.6em}{}

\titlespacing*{\section}{0pt}{0.95em}{0.35em}
\titlespacing*{\subsection}{0pt}{0.65em}{0.25em}
\titlespacing*{\subsubsection}{0pt}{0.55em}{0.2em}

% ---------- Table engine ----------
\usepackage{tabularray}
\UseTblrLibrary{booktabs}

% ---------- tabularray: GLOBAL table treatment ----------
\SetTblrInner{
  cells = {fg=black, bg=white, valign=t},
}

% ---------- Header helpers ----------
\newcommand{\Hdr}[1]{\shortstack[c]{\strut #1\strut}}

% ---------- Lines ----------
\newcommand{\TitleLine}{\textcolor{StageBlue}{\rule{0.98\linewidth}{0.95pt}}}
\newcommand{\SubTitleLine}{\textcolor{RuleGray}{\rule{0.98\linewidth}{0.65pt}}}

% ---------- Continued on next page ----------
\newcommand{\ContinuedOnNextPageInline}{%
  \par
  \vfill
  \noindent\textcolor{RuleGray}{\rule{\linewidth}{1pt}}\vspace{-2pt}
  \noindent\hfill\textit{Continued on next page}\par\vspace{-10pt}
  \noindent\textcolor{RuleGray}{\rule{\linewidth}{1pt}}
  \clearpage
}

% ---------- PDF hyperlinks ----------
\usepackage[hidelinks]{hyperref}
\hypersetup{
  colorlinks=false,
  pdfborder={0 0 0},
  linktoc=all,
  pdftitle={\DocTitle},
  pdfauthor={\ProgramName}
}

% =========================================================
\begin{document}

\section*{7.11 — Footwear Systems and Foot Care}
\phantomsection
\addcontentsline{toc}{section}{7.11 — Footwear Systems and Foot Care}

Overview \textbf{ID}: 7.11\\
Overview Name: Footwear Systems and Foot Care\\
Version: v1.0\\
Governing Status: Draft — Non-Deployable\\
Scope Boundary Statement: Covers footwear selection, sock interfaces, blister prevention posture, lacing and fit checks, drying and storage posture, and non-medical foot-care controls; does not provide workouts, protocols, medical guidance, or load-carriage system prescriptions.

\subsection*{7.11.1 Purpose \& Scope}
\phantomsection
\addcontentsline{toc}{subsection}{7.11.1 Purpose \& Scope}

This category covers:

\begin{itemize}
\item Footwear selection by modality: walking, running, rucking, general training.
\item Sock systems and blister-prevention interfaces (liner/outer systems, seam/fit discipline).
\item Lacing/fit checks, replacement cues, storage/drying posture, and travel continuity.
\item Foot-care consumables and routines that stabilize adherence, reduce injury risk, and reduce validity drift for route repeatability.
\end{itemize}

This overview crosswalks Stage 6 timing to the footwear and foot-care conditions required for safe execution and admissible evidence.

This category does not cover:

\begin{itemize}
\item Load carriage systems: Stage 7.07.
\item Apparel layers beyond socks: Stage 7.14 and Stage 7.15.
\item Hygiene / sanitation beyond foot-related items: cross-reference Stage 7.20.
\item Testing / timing tools: cross-reference Stage 7.18.
\end{itemize}

This overview does not provide programming, test protocols, calendars, medical diagnosis, or treatment guidance.

\subsection*{7.11.2 Governance And Safety Boundaries}
\phantomsection
\addcontentsline{toc}{subsection}{7.11.2 Governance And Safety Boundaries}

\begin{itemize}
\item Stage 0 boundary alignment: stop rules, safety override doctrine, conflict-resolution hierarchy, and evidence-admissibility controls apply to all 7.11 selections and substitutions.
\item Footwear functions as safety infrastructure and a validity stabilizer. Poor fit increases injury risk and increases invalid testing risk by changing gait mechanics and comparability.
\item Start-from-zero posture: Phase 00 requires a minimum viable, stable footwear interface aligned to Phase 00 low-impact bias.
\item No push-through rule:
  \begin{itemize}
  \item gait-altering pain stops the exposure today,
  \item substitute lower-risk modality,
  \item document per validity/logging posture.
  \end{itemize}
\item Escalation triggers requiring stop and seek evaluation (non-medical):
  \begin{itemize}
  \item rapidly worsening swelling/redness/heat,
  \item drainage from a blister or wound site,
  \item fever or systemic illness signs,
  \item numbness or tingling,
  \item gait-altering pain.
  \end{itemize}
\item Evidence inadmissibility: testing or gate evidence produced in clearly unstable, unsafe, or unvalidated footwear is inadmissible until corrected and re-verified.
\item Substitution posture: when a planned footwear interface is unsafe or unvalidated, substitute to the lowest-risk modality that preserves intent rather than forcing exposure.
\item “No novelty near Deload+Test” posture applies:
  \begin{itemize}
  \item do not first-use new shoes/socks, new insoles, new traction aids, or a new lacing system inside protected windows.
  \end{itemize}
\end{itemize}

\subsection*{7.11.3 Tier Doctrine}
\phantomsection
\addcontentsline{toc}{subsection}{7.11.3 Tier Doctrine}

\begin{itemize}
\item Price band convention: \$ = minimal household-cost item; \$\$ = modest purpose-built item; \$\$\$ = expanded-capability or redundancy item.
\item N/A rule: use N/A only when a field is genuinely not applicable and omission does not obscure fit, maintenance, or phase-timing logic.
\item Substitution posture: substitutions are acceptable only when they preserve safety, validated fit, and movement intent; substitutes that introduce gait instability, blister escalation, or invalid test conditions break intent.
\item Tier 0: household-only / inadequate; temporary constraint state only; not deployable for Phase 00 (End) execution.
\item Tier 1: minimum viable safe footwear interface + socks + blister kit + nail care + drying posture to support Phase 00 (End) and early progression.
\item Tier 2: expanded modality-specific rotation and environmental control:
  \begin{itemize}
  \item traction aids (ice control),
  \item expanded blister prevention supplies,
  \item OTC comfort insoles (comfort-only posture),
  \item expanded sock options and redundancy.
  \end{itemize}
\item Tier 3: advanced and overbuilt redundancy and travel continuity:
  \begin{itemize}
  \item multiple pairs staged per modality,
  \item advanced blister prevention redundancy,
  \item clinician-prescribed orthotics posture (clinician-directed only; not required; no medical guidance provided).
  \end{itemize}
\end{itemize}

\subsection*{7.11.4 Selection Criteria \& Fit Checks}
\phantomsection
\addcontentsline{toc}{subsection}{7.11.4 Selection Criteria \& Fit Checks}

Fit and safety checklist (operational):

\begin{itemize}
\item Toe box clearance: adequate length and width to protect nails and accommodate swelling.
\item Heel lockdown: no heel slip during walking and under load; adjust lacing before changing shoes.
\item Midfoot stability: stable platform under lateral movement and load.
\item Modality specificity (non-programming posture):
  \begin{itemize}
  \item walking shoe prioritizes comfort/stability for long low-impact duration,
  \item Run shoe prioritizes repeatable gait pattern and cushioning consistency,
  \item Ruck-compatible footwear prioritizes stability under load and sock-volume accommodation.
  \end{itemize}
\item Sock interface fit:
  \begin{itemize}
  \item avoid seam pressure points,
  \item ensure shoe volume accommodates liner+outer system when required.
  \end{itemize}
\item Wear-test rule:
  \begin{itemize}
  \item validate any new footwear/sock combination on short, low-risk sessions first,
  \item do not first-use new footwear on Deload+Test windows or test days.
  \end{itemize}
\end{itemize}

\subsection*{7.11.5 Setup, Safety, and Anchoring}
\phantomsection
\addcontentsline{toc}{subsection}{7.11.5 Setup, Safety, and Anchoring}

Setup doctrine (no programming):

\begin{itemize}
\item Lacing protocol (intent-level):
  \begin{itemize}
  \item heel-lock posture if needed,
  \item pressure point checks and re-lace rather than “accept discomfort,”
  \item maintain consistent lacing approach during decision windows.
  \end{itemize}
\item Moisture management:
  \begin{itemize}
  \item dry shoes after every session; rotate pairs when possible.
  \end{itemize}
\item Surface safety:
  \begin{itemize}
  \item wet/ice traction is a safety constraint; traction aids are Tier 2 / Tier 3 only.
  \item when unsafe, substitute indoors rather than attempting “careful walking.”
  \end{itemize}
\item Change-control posture:
  \begin{itemize}
  \item avoid changing footwear system and introducing new ruck interface variables in the same week,
  \item log footwear changes and hotspot events immediately.
  \end{itemize}
\end{itemize}

\subsection*{7.11.6 Maintenance, Replacement, and Storage}
\phantomsection
\addcontentsline{toc}{subsection}{7.11.6 Maintenance, Replacement, and Storage}

\begin{itemize}
\item Replacement cues:
  \begin{itemize}
  \item outsole traction loss,
  \item cushioning collapse,
  \item new persistent hotspots that were not present previously,
  \item heel counter breakdown or loss of midfoot stability.
  \end{itemize}
\item Consumable cadence:
  \begin{itemize}
  \item sock rotation and replacement (replace thin/holes early),
  \item blister supply restock (tape/pads).
  \end{itemize}
\item Storage:
  \begin{itemize}
  \item dry and ventilated; avoid leaving shoes wet in closed bags.
  \item staged pairs kept ready to prevent last-minute substitutions.
  \end{itemize}
\end{itemize}

\subsection*{7.11.7 Substitution Rules — Maintain Intent Without Breaking Safety and Validity}
\phantomsection
\addcontentsline{toc}{subsection}{7.11.7 Substitution Rules — Maintain Intent Without Breaking Safety and Validity}

\begin{itemize}
\item If footwear is inadequate for planned exposure:
  \begin{itemize}
  \item substitute to the lowest-risk modality that preserves cardio intent (e.g., indoor low-impact tool from 7.08),
  \item do not perform running or load-bearing exposures in unstable footwear,
  \item document the deviation.
  \end{itemize}
\item Prohibited substitutions:
  \begin{itemize}
  \item running in poor-fit or unstable footwear,
  \item load-bearing exposures when footwear interface is unvalidated,
  \item “new shoe on test day” behavior.
  \end{itemize}
\item Stage 2.E validity/logging fields conceptually required when substitutions occur:
  \begin{itemize}
  \item Route \textbf{ID} / surface condition notes (or indoor modality identity),
  \item footwear used and its condition,
  \item reason-code posture when a deviation occurs,
  \item validity impact posture when comparability is compromised (reslot rather than claim gate credit).
  \end{itemize}
\end{itemize}

\subsection*{7.11.8 Phase Integration Map — Required}
\phantomsection
\addcontentsline{toc}{subsection}{7.11.8 Phase Integration Map — Required}

\begingroup
\scriptsize
\setlength{\tabcolsep}{2pt}
\renewcommand{\arraystretch}{1.12}

\begin{longtblr}{
  width=\linewidth,
  colspec = {
    Q[c,wd=1.05in]
    Q[c,wd=1.55in]
    X[c,3.05]
    X[c,3.95]
  },
  rowhead=1,
  hlines,
  vlines,
  cells = {hsep=2pt, vsep=6pt, halign=c, valign=m},
  row{1} = {bg=StageBlue!15, fg=black, font=\bfseries, valign=m},
  row{odd} = {bg=white},
  row{even} = {bg=LightGray},
}
\Hdr{Phase} &
\Hdr{Required Tier (from Stage 6)} &
\Hdr{What Changes / Becomes Mandatory} &
\Hdr{Notes} \\
Phase 00 &
Tier 1 &
Minimum viable footwear interface becomes mandatory: supportive walking/training shoe + baseline sock rotation + basic blister kit + nail care &
Phase 00 (End) is the first required point. Do not change shoe model/class inside protected windows without stabilization. \\
Phase 01 &
Tier 1 &
Maintain stable footwear interface; introduce/lock a run-appropriate shoe option if \textbf{RUN} evidence becomes phase-relevant &
Phase 01 (End): treat run shoe choice as a controlled variable during decision windows. \\
Phase 02 &
Tier 1 &
Maintain stable walking/run footwear; increase emphasis on foot integrity monitoring as training consequence rises &
Phase 02 (End): do not let strength buildout drive foot-neglect drift. \\
Phase 03 &
Tier 1 &
Sock system milestone: liner+outer becomes mandatory by Phase 03 (End) aligned to first ruck Tier 1 milestone (Stage 7.07) &
Phase 03 (End) is the first ruck Tier 1 timing. Stabilize liner+outer system before protected windows; do not first-use on gate attempts. \\
Phase 04 &
Tier 1 &
Maintain liner+outer sock system; ensure ruck-compatible footwear interface is stable for Phase 04 (End) ruck feasibility &
Phase 04 (End) ruck feasibility increases consequence. Even though overall category tier remains Tier 1, the ruck footwear interface must be stabilized (no novelty near protected windows). \\
Phase 05 &
Tier 1 &
Maintain stable interfaces; longer endurance increases cost of hotspots and toenail trauma &
Phase 05 (End): replace worn socks/shoes early rather than pushing through. \\
Phase 06 &
Tier 2 &
Tier 2 becomes mandatory: traction aids readiness, expanded blister prevention redundancy, OTC comfort insoles optional &
Phase 06 (End) is the first Tier 2 requirement. Ice/traction becomes a continuity constraint; traction aids are Tier 2 / Tier 3 only. \\
Phase 07 &
Tier 2 &
Maintain Tier 2; enforce stricter no-novelty posture due to hybrid density &
Phase 07 (End): prevent interface changes that increase variance and injury risk. \\
Phase 08 &
Tier 2 &
Maintain Tier 2; load tolerance increases foot breakdown risk; blister prevention is a primary control &
Phase 08 (End): foot risk becomes a dominant gate constraint; substitute/reslot rather than forcing. \\
Phase 09 &
Tier 2 &
Maintain Tier 2; high \textbf{CAL} volume increases friction and skin risk; maintain redundancy &
Phase 09 (End): prevent skin breakdown that drives gait drift. \\
Phase 10 &
Tier 2 &
Maintain Tier 2; multi-domain gates increase consequence of footwear variance &
Phase 10 (End): treat footwear changes as decision-window-sensitive variables. \\
Phase 11 &
Tier 2 &
Maintain Tier 2; recovery sensitivity increases consequence of foot breakdown &
Phase 11 (End): prioritize stability and early replacement. \\
Phase 12 &
Tier 3 &
Tier 3 becomes mandatory: modality-staged redundancy, travel kit redundancy, advanced blister prevention redundancy &
Phase 12 (End): redundancy is introduced to reduce risk of compromised test windows due to footwear failure or last-minute changes. \\
Phase 13 &
Tier 3 &
Maintain Tier 3; consolidation demands low variance and stable interfaces &
Phase 13 (End): no novelty near final protected windows; repeatability posture dominates. \\
\end{longtblr}

\endgroup

7.11.8 Phase Integration Map Notes:

\begin{itemize}
\item Sock liner+outer system is explicitly tied to the first ruck Tier 1 milestone (Phase 03 (End)). If any upstream artifact schedules ruck exposure earlier, Requires Stage 8 review.
\end{itemize}

\subsection*{7.11.9 Risks, Contraindications, and Red Flags}
\phantomsection
\addcontentsline{toc}{subsection}{7.11.9 Risks, Contraindications, and Red Flags}

\begin{itemize}
\item Hotspots/blisters leading to gait alteration:
  \begin{itemize}
  \item hotspots must be treated immediately; blister progression is a gate-threatening failure mode.
  \end{itemize}
\item Toenail trauma:
  \begin{itemize}
  \item poor toe box clearance and downhill friction increase risk; manage fit and replace shoes when toe trauma trends appear.
  \end{itemize}
\item Plantar pain under load:
  \begin{itemize}
  \item treat as a stop-and-downshift signal; do not force ruck/run exposures through escalating pain.
  \end{itemize}
\item Falls on slick surfaces:
  \begin{itemize}
  \item traction aids are Tier 2 / Tier 3 safety controls; if still unsafe, substitute indoors/reslot.
  \end{itemize}
\item Red flags requiring stop and seek evaluation (non-medical):
  \begin{itemize}
  \item gait-altering pain,
  \item numbness or tingling,
  \item rapidly worsening swelling/redness/heat,
  \item drainage from a blister or wound site,
  \item fever or systemic illness signs.
  \end{itemize}
\item Validity risk:
  \begin{itemize}
  \item unlogged footwear changes can contaminate comparability; document footwear used and condition; reslot gate claims when novelty occurs.
  \end{itemize}
\item Timing mismatch trigger:
  \begin{itemize}
  \item if any phase artifact demands footwear Tier 2/3 earlier than Stage 6, or demands ruck footwear before ruck Tier 1 timing, Requires Stage 8 review.
  \end{itemize}
\end{itemize}

\end{document}